\section{Exact outgoing-star deletion constraints in the original graph}
\label{sec:outstar-deletion-criticality}

For an original oriented graph $D$ and a vertex $f$, put
$O=N_D^+(f)$, $B=N_D^{++}(f)$, and $g=g_D(f)$.
For every $b\in B$ let
\[
 I_b=N_D^-(b)\cap O
\]
be its nonempty set of original first-neighborhood witnesses.
For $\varnothing\ne S\subseteq O$, delete precisely the original
arcs $f\to S$, obtaining $D_S$, and define
\[
 L(S)=\{b\in B:I_b\subseteq S\},\qquad
 R(S)=S\cap\Gamma_D(O\setminus S).
\]
Thus $L(S)$ is the set of lost old second targets, while $R(S)$ is
the set of removed first targets which become exact second targets.

\begin{lemma}[Exact loss and retargeting formula]
\label{lem:outstar-exact-loss-retargeting}
For every original oriented $D$,
\[
 N_{D_S}^{++}(f)=(B\setminus L(S))\,\dot\cup\,R(S),
\]
and hence
\[
 g_{D_S}(f)=g-|S|+|L(S)|-|R(S)|.
\]
If $D$ is all-deficient and minimal under deletion of arbitrary
nonempty arc sets on its fixed vertex set, then
\[
 |L(S)|\le |S|+|R(S)|-g
 \quad\text{for every }\varnothing\ne S\subseteq O.
 \tag{OD1}\label{eq:outstar-minimal-budget}
\]
\end{lemma}
\begin{proof}
The retained first set is $O\setminus S$. An old second target
$b$ remains reached exactly when $I_b$ meets that retained set.
No originally far target can become reached when arcs are deleted.
The only possible new exact second targets are removed first targets
in $S$, and precisely those reached from a retained intermediate form
$R(S)$. These two classes are disjoint, proving the exact identity.

Every vertex other than $f$ has unchanged first neighborhood and can
only lose second targets. Thus any Seymour vertex forced by minimality
of the nonempty deletion must be $f$, giving $g_{D_S}(f)\le0$.
Substitute the exact gap formula.
\end{proof}

\begin{corollary}[Every proper retained set is externally expanding]
\label{cor:outstar-proper-retained-expansion}
In the same minimal class, every proper subset $X\subsetneq O$ satisfies
\[
 |\Gamma_D(X)\setminus X|\ge|X|,
 \qquad |\Gamma_D(O)\setminus O|=|O|-g<|O|.
\]
Thus $O$ is inclusion-minimal among its subsets that fail this
external expansion inequality.
\end{corollary}
\begin{proof}
For $S=O\setminus X$, the retained first set is $X$. No member of $X$
points to $f$, because each original $f\to x$ excludes its reverse.
Therefore $N_{D_S}^{++}(f)=\Gamma_D(X)\setminus X$.
Minimality gives the first inequality. With $X=O$, the same external
boundary is precisely $B$, giving the strict deficit. The empty
retained set has both sides zero and causes no exception.
\end{proof}

\begin{corollary}[Gap two has no private old second heads]
\label{cor:outstar-gap-two-multiple-witnesses}
If $g_D(f)=2$ in the original minimal class, then for every $a\in O$:
\begin{enumerate}
\item Some other original first vertex points to $a$.
\item No old second target has $a$ as its unique original first witness.
\item Deleting just $f\to a$ preserves all old exact second targets and
adds precisely $a$ as an exact second target. The new gap at $f$ is zero.
\end{enumerate}
Consequently
\[
 \delta^-(D[O])\ge1,\qquad
 |I_b|\ge2\quad(b\in B).
\]
The induced original first-neighborhood graph contains a directed cycle
of length at least three. For $g_D(f)=1$, the singleton constraint is
instead $|L(\{a\})|\le|R(\{a\})|\le1$; it does not force a cycle.
\end{corollary}
\begin{proof}
For $S=\{a\}$, the retargeting set has size zero or one and the loss
set has nonnegative size. At gap two, \eqref{eq:outstar-minimal-budget}
gives $0\le|L(S)|\le|R(S)|-1\le0$. Thus $|R(S)|=1$ and $L(S)$
is empty. The exact identities prove all three assertions. A head with
a unique first witness would belong to that singleton loss set, so every
$I_b$ has at least two vertices. Following predecessors in $D[O]$ produces
a cycle; orientedness excludes lengths one and two. The gap-one bound
is the same formula with $g=1$.
\end{proof}

\paragraph{Use at closure feeds.}
These inequalities apply to the original first set of every final
closure feed, since that set has been proved unchanged. They constrain
original two-step witnesses, not new second heads added by closure or
completion. No intermediate closure graph is assumed arc-minimal, and
deleting its newly added arcs is not justified by original minimality.
The conditions are necessary, not sufficient to construct a graph with
all vertices deficient.
